\section{Evaluation and real-data transfer}
\label{sec:rq3}
\label{sec:lidar-evaluation}

Evaluation must distinguish similarity between scans from benefits obtained
by training on them. The reviewed studies provide evidence at both levels,
but their different tasks and protocols require separate interpretation.

\subsection{Scan diagnostics and training evidence}

Geometric errors require aligned, corresponding observations, whereas
unpaired generation needs distributional comparisons. LiDARGen uses
\gls{mmd} and \gls{jsd} on bird's-eye-view (\gls{bev}) histograms and
\gls{frd} on RangeNet++ features. On KITTI-360, ProjectedGAN obtains better
MMD, while LiDARGen obtains better FRD and JSD
\parencite{zyrianovLearningGenerateRealistic2022}.
The disagreement shows that conclusions depend on the representation used
for comparison. These metrics diagnose differences but do not establish
which generated dataset provides better training data.

Training protocols provide complementary evidence. PreSIL reports improved
KITTI pedestrian detection after synthetic pretraining and real
\gls{fine-tuning}. However, selecting its best checkpoint on the evaluation
split can produce optimistic results relative to independent final testing
\parencite{hurlPreciseSyntheticImage2019}.
LiDARsim reports a KITTI hard BEV score of 73.7 at \gls{iou}~0.7, compared
with 42.2 for modified CARLA and 80.0 for real training
\parencite{manivasagamLiDARsimRealisticLiDAR2020}.
This narrows the simulation-to-real difference in that setting, but uses
sensor adaptation and a KITTI-labelled object bank. It is consequently
neither evidence of independence from real data nor a full 3D detection result.

SynLiDAR's mixed-training results (Table~\ref{tab:transfer}) show gains of
2.2 and 2.6 percentage points over real-only training. PCT increases those
gains to 4.4 and 5.2 points \parencite{xiaoTransferLearningSynthetic2022}.
Thus, supplementation helps in the reported settings, while the additional
improvement with translation cannot be attributed to sample quantity alone.
LiDARGen instead applies an existing segmenter to densified scans without
retraining, supporting completion utility rather than synthetic-training
gains \parencite{zyrianovLearningGenerateRealistic2022}.

\begin{table}[htbp]
\caption{SynLiDAR mIoU (\%) reported by Xiao et al., Tables~2--3
\parencite{xiaoTransferLearningSynthetic2022}.
All rows use MinkowskiNet. Evaluation uses SemanticKITTI validation
sequence 08 and SemanticPOSS validation sequence 03.}
\label{tab:transfer}
\small
\begin{tabularx}{\linewidth}{@{}Yrr@{}}
\toprule
\textbf{Training data} & \textbf{SemanticKITTI} & \textbf{SemanticPOSS} \\
\midrule
Real only & 60.3 & 50.6 \\
Real + SynLiDAR & 62.5 & 53.2 \\
Real + PCT-translated SynLiDAR & 64.7 & 55.8 \\
\bottomrule
\end{tabularx}
\end{table}

\subsection{Validity and comparison conditions}

Validation supports model selection, while independent real test data
support final performance estimates
\parencite{raschkaModelEvaluationModel2020}.
Westphal and Brannath address selection-induced optimism with a specialised
multiple-testing correction, which does not justify ordinary repeated
test-set tuning \parencite{westphalImprovingModelSelection2019}.
For LiDAR, independence must also extend to generator inputs and
reconstructed assets. Comparisons should document scene splits, sensor
settings, class mappings, training budgets and access to real labels.

Task-appropriate real-data \gls{ap} or \gls{miou}, supported by repeated
runs and class- and range-specific analysis, then provides direct evidence
of training utility. Sensor diagnostics help explain the result.
Evaluation of an already-trained autonomy system is a different objective:
aggregate perception scores need not track planning discrepancies
\parencite{manivasagamZeroDomainGap2023}.
These distinctions determine how broadly the reviewed findings can be interpreted.
